\section{Moving-target local saturation}\label{sec:moving-target}

Keep the fixed reduced target $t=a/b\in(0,1)$, exponent $\gamma$, direct family $\mathcal E=\mathcal E_X$, common denominator $L=L_X$, and Bernoulli parameter $\theta=t/\Lambda$. For a lattice target $s\in L^{-1}\Z$, write
\[
W_X(s)=
\Pr\left(\sum_{e\in\mathcal E}\frac{\xi_e}{e}=s\right).
\tag{8.1}\label{eq:8.1}
\]
We prove the full statement of Theorem \ref{thm:local-profile}.

\subsection{The target shift is a unit character}

Let
\[
t_j=t+\frac{j}{L},
\qquad j\in\Z.
\tag{8.2}\label{eq:8.2}
\]
The quotient coefficient at $t_j$ is
\[
L W_X^{\mathrm{quot}}(t_j)
=
\sum_{h\bmod L}F_j(h),
\tag{8.3}\label{eq:8.3}
\]
where
\[
F_j(h)
=
e(-ht_j)
\prod_{e\in\mathcal E}
\bigl((1-\theta)+\theta e(h/e)\bigr).
\]
Comparing with the central target $t_0=t$ gives the exact identity
\[
F_j(h)=e(-hj/L)F_0(h).
\tag{8.4}\label{eq:8.4}\label{eq:unit-character}
\]
Hence
\[
|F_j(h)|=|F_0(h)|
\tag{8.5}\label{eq:8.5}
\]
for every $h$ and $j$. All anchor, row, decoder, compression, and coherent-tail estimates in Sections \ref{sec:anchor}--\ref{sec:gaussian} are therefore unchanged.

\subsection{The Gaussian phase}

Fix $U>0$ and consider integers satisfying
\[
|j|\le UL\sigma.
\tag{8.6}\label{eq:8.6}
\]
Put
\[
u_j=\frac{j}{L\sigma},
\qquad |u_j|\le U.
\tag{8.7}\label{eq:8.7}
\]
On a genuine integer frequency $m$, equations \eqref{eq:6.10} and \eqref{eq:8.4} give, uniformly for $|m|\le C/\sigma$,
\[
F_j(m)
=
\exp\left(
-2\pi^2m^2\sigma^2
-\frac{2\pi i mj}{L}
+\varepsilon_{X,m}
\right).
\tag{8.8}\label{eq:8.8}
\]
With $x_m=m\sigma$, the new phase is
\[
\exp(-2\pi i x_m u_j).
\tag{8.9}\label{eq:8.9}
\]

\begin{proposition}[Compact-uniform major limit]\label{prop:moving-major}
For fixed $C,U>0$, uniformly for integers $j$ satisfying \eqref{eq:8.6},
\[
\sigma
\sum_{|m|\le C/\sigma}F_j(m)
-
\int_{-C}^{C}e^{-2\pi^2x^2}e(-u_jx)\,dx
\longrightarrow0.
\tag{8.10}\label{eq:8.10}
\]
\end{proposition}

\begin{proof}
The Taylor remainder is the same uniform error as in \eqref{eq:6.10}. For $|u|\le U$, the functions
\[
x\mapsto e^{-2\pi^2x^2}e(-ux)
\]
have uniformly bounded derivative on $[-C,C]$. The Riemann-sum approximation is therefore uniform in $u\in[-U,U]$.
\end{proof}

The full Fourier transform is
\[
\int_{\R}e^{-2\pi^2x^2}e(-ux)\,dx
=
\Phi(u)
:=
\frac1{\sqrt{2\pi}}e^{-u^2/2}.
\tag{8.11}\label{eq:8.11}
\]
The truncation tail outside $[-C,C]$ is bounded by the absolute Gaussian tail, uniformly in $u$.

\subsection{Uniform complementary estimates}

By \eqref{eq:8.5}, every modulus estimate in Sections 6.2--6.5 is unchanged. If $\mathfrak R_{X,C}(j)$ is the complete contribution outside the genuine integer range $|m|\le C/\sigma$, then
\[
\limsup_{X\to\infty}
\sup_{|j|\le UL\sigma}
\sigma|\mathfrak R_{X,C}(j)|
\le
K e^{-cC^2/2}.
\tag{8.12}\label{eq:8.12}
\]
No factor counting the number of targets appears; the estimate is already uniform in the inserted unit character.

Combining Proposition \ref{prop:moving-major} with \eqref{eq:8.12}, first taking $X\to\infty$ with $C,U$ fixed and then $C\to\infty$, gives
\[
L\sigma W_X^{\mathrm{quot}}(t_j)
=
\Phi(u_j)+o_U(1)
\tag{8.13}\label{eq:8.13}
\]
uniformly for \eqref{eq:8.6}.

\subsection{No-wrap and simultaneous exact realization}

By \eqref{eq:2.9} and \eqref{eq:2.10}, the fixed centre $t$ lies a positive distance from both endpoints of the eventual interval $[0,\Lambda]$. Since $\sigma\to0$, for each fixed $U$ all targets $t_j$ with \eqref{eq:8.6} lie in $[0,\Lambda]$ for sufficiently large $X$. Lemma \ref{lem:no-wrap} therefore gives
\[
W_X^{\mathrm{quot}}(t_j)=W_X(t_j).
\tag{8.14}\label{eq:8.14}
\]
Together with \eqref{eq:8.13}, this proves
\[
L\sigma
W_X\!\left(t+\frac{j}{L}\right)
=
\Phi\!\left(\frac{j}{L\sigma}\right)+o_U(1)
\tag{8.15}\label{eq:8.15}\label{eq:moving-profile}
\]
uniformly for $|j|\le UL\sigma$, which is Theorem \ref{thm:local-profile}.

Since $\Phi$ has a positive minimum on $[-U,U]$, every coefficient in this window is eventually positive simultaneously. Hence the integer subset-sum image
\[
\mathcal S_X
=
\left\{\sum_{e\in A}\frac{L}{e}:A\subseteq\mathcal E\right\}
\tag{8.16}\label{eq:8.16}
\]
contains the entire block
\[
\left\{Lt+j:
|j|\le\lfloor UL\sigma\rfloor\right\}.
\tag{8.17}\label{eq:8.17}
\]
By \eqref{eq:2.19} and Lemma \ref{lem:family-size},
\[
\sigma\to0,
\qquad
L\sigma\to\infty.
\tag{8.18}\label{eq:8.18}
\]
Thus the number of represented lattice points tends to infinity while their radius after division by $L$ is at most $U\sigma\to0$.

Dividing \eqref{eq:8.15} by the central local limit also gives
\[
\frac{W_X(t+j/L)}{W_X(t)}
=
\exp\left(-\frac{j^2}{2L^2\sigma^2}\right)+o_U(1)
\tag{8.19}\label{eq:8.19}
\]
uniformly on the same window.

\begin{remark}
There is no largest finite admissible $U$: every fixed $U$ works after a threshold depending on $U$. The proof does not give a growing standardized radius $U=U_X\to\infty$; such a statement would require quantitative errors small enough to compete with $e^{-U_X^2/2}$.
\end{remark}
